% Carl Friedrich Gauss, Werke IX, printed 304–306, Article [5] through case I.
\par\bigskip
\begin{center}
[5.]\\[4pt]
Bedingungsgleichungen der zweiten Art.
\end{center}

[Nach der Ausgleichung der Winkelschlussfehler der Dreiecke werden die
Bedingungsgleichungen aufgestellt, die sich aus den Seitenverhältnissen ergeben.

Für das Viereck $7.\ 8.\ 10.\,.\,.\,.\,.\,.\ 9$ gilt bei Wahl von 9 als Zentralpunkt:
\begin{center}
\small
\setlength{\tabcolsep}{3pt}
\renewcommand{\arraystretch}{1.25}
\begin{tabular}{@{}r r@{: }r r r l@{}}
\shortstack{im\\Dreieck}&\multicolumn{2}{c}{ausgegl. Winkel}&$\frac13$ Exzess&$\log\sin$&Verbesserung\quad]\\
7 auf&$7$&$66^\circ\ 1'\ 19''{,}151$&$-2''{,}269$&$9{,}960\,8021.983$&$+9.36\{(7.9)-(7.8)\}$\\
&$8$&$66\ \ 39\ \ 58{,}666$&$-2{,}269$&$9{,}962\,9416.342$&$+9.08\{(8.7)-(8.9)\}$\\[4pt]
$9$&$10$&$22\ \ 59\ \ 17{,}428$&$-0{,}794$&$9{,}591\,6628.409$&$+49.63\{(10.7)-(10.9)\}$\\
&$7$&$10\ \ 27\ \ 3{,}217$&$-0{,}794$&$9{,}258\,6108.427$&$+114.15\{(7.9)-(7.10)\}$\\
\end{tabular}
\end{center}
\newpage
\noindent\textbf{ZUR NETZAUSGLEICHUNG.}\hfill305
\begin{center}
\small
\setlength{\tabcolsep}{3pt}
\renewcommand{\arraystretch}{1.25}
\begin{tabular}{@{}r r@{: }r r r l@{}}
$\bigl[$\shortstack{im\\Dreieck}&\multicolumn{2}{c}{ausgegl. Winkel}&$\frac13$ Exzess&$\log\sin$&Verbesserung\quad]\\
10 auf&$8$&$23^\circ\ 11'\ 52''{,}167$&$-1''{,}415$&$9{,}595\,3867.516$&$+49.13\{(8.9)-(8.10)\}$\\
&$10$&$57\ \ 33\ \ 19{,}333$&$-1{,}415$&$9{,}926\,2943.910$&$+13.39\{(10.8)-(10.9)\}$\\
&&&&$9{,}997\,8605.641$&\\
&&&&$0{,}333\,0519.982$&\\
&&&&$9{,}669\,0923.606$&\\
\cline{5-5}
&&&&$+49.229$&\\
&&&&\textnormal{(vorher $+57$) [*)]}&\\
\end{tabular}
\end{center}

[Die Seitengleichung für $7.\ 8.\ 10.\,.\,.\,.\,.\,.\ 9$ lautet daher:]
\[
\begin{aligned}
0={}&+49{,}229-9{,}36(7.8)-104{,}79(7.9)+114{,}15(7.10)\\
&-9{,}08(8.7)+58{,}21(8.9)-49{,}13(8.10)\\
&+49{,}63(10.7)-13{,}39(10.8)-36{,}24(10.9).
\end{aligned}
\]

[Diese Gleichung wird wie folgt umgeformt. Ist bei der Ausgleichung
einer Figur neben einer Anzahl von Winkelgleichungen, deren konstante Glieder
null sind, nur eine Seitengleichung zu erfüllen, so genügt die Ausgleichung
\textls[120]{einer einzigen} Gleichung. Man erhält diese Gleichung, indem man
die Winkelgleichungen mit bestimmten Faktoren multipliziert
und zur Seitengleichung addiert. Die Faktoren sind so zu bestimmen,
dass die Summe der Quadrate der Koeffizienten in der umgeformten
Seitengleichung minimal wird. Man nehme also die bereits ausgeglichenen
Winkelgleichungen der Dreiecke 7, 9, 10:
\[
\begin{aligned}
+(7.9)-(9.7)+(9.8)-(8.9)+(8.7)-(7.8)&=0\\
+(7.9)-(9.7)+(9.10)-(10.9)+(10.7)-(7.10)&=0\\
+(8.9)-(9.8)+(9.10)-(10.9)+(10.8)-(8.10)&=0,
\end{aligned}
\]
Multipliziert man sie mit $x,y,z$ und addiert sie zur obigen Seitengleichung,
so ergibt sich zunächst:
\[
\begin{aligned}
0={}&+49{,}229-(9{,}36+x)(7.8)-(104{,}79-x-y)(7.9)+(114{,}15-y)(7.10)\\
&-(9{,}08-x)(8.7)+(58{,}21-x+z)(8.9)-(49{,}13+z)(8.10)\\
&-(x+y)(9.7)+(x-z)(9.8)+(y+z)(9.10)\\
&+(49{,}63+y)(10.7)-(13{,}39-z)(10.8)-(36{,}24+y+z)(10.9),
\end{aligned}
\]
\par\smallskip
\noindent\rule{0.25\textwidth}{0.4pt}\par
{\footnotesize
[*) Der eingeklammerte Wert $+57$ ergibt sich, wenn die Seitengleichung mit den beobachteten
Winkelwerten aus Art. 3 berechnet wird.]\par}
\vfill
\noindent\textsc{IX.}\hfill39
\newpage
\noindent306\hfill\textbf{NACHLASS.}\par\medskip
\noindent und minimiert man nun die Summe der Quadrate der Koeffizienten
der Verbesserungen, so müssen
\[
\begin{aligned}
6x+2y-2z&=+162{,}7\\
2x+6y+2z&=+133{,}1\\
-2x+2y+6z&=-130{,}2
\end{aligned}
\]
gelten.] Man erhält
\[
\begin{aligned}
x&=+7{,}76\\
y&=+29{,}21\\
z&=-28{,}85,
\end{aligned}
\]
und damit die so verbesserte Bedingungsgleichung:
\[
\begin{aligned}
\mathrm{I.}\quad 0={}&+49{,}229-17{,}12(7.8)-67{,}82(7.9)+84{,}94(7.10)\\
&-1{,}32(8.7)+21{,}60(8.9)-20{,}28(8.10)\\
&-36{,}97(9.7)+36{,}61(9.8)+0{,}36(9.10)\\
&+78{,}84(10.7)-42{,}24(10.8)-36{,}60(10.9).
\end{aligned}
\]

[Im Folgenden sind nur die ursprünglichen und die umgeformten
Seitengleichungen angegeben. Die Überschrift nennt jeweils die Figur, zu der
die betreffende Seitengleichung gehört, und die Dreiecke, deren Winkelgleichungen
zu ihrer Umformung verwendet werden. Anstelle der Faktoren
$x,y,z$ stehen die Dreiecksnummern.]
